\begin{essay}{}
	Расчетно-пояснительная записка содержит \pageref{LastPage} страниц, 4 раздела, 7 рисунков, 10 таблиц, 19 использованных источников, 1 приложение.

	В работе представлена разработка метода обеспечения согласованности состояний двух асинхронных игроков с помощью конечных автоматов. Проведен анализ подходов к синхронизации игровых состояний и передаче данных в многопользовательских приложениях, рассмотрены причины рассинхронизации при UDP-сетевом обмене. Разработан метод событийной синхронизации, включающий подтверждение обработки событий, повторную отправку неподтвержденных сообщений, контроль порядка и фильтрацию дубликатов. Выполнена программная реализация метода и проведено исследование его эффективности при задержках, потерях и дублировании UDP-пакетов. Проведено сравнение разработанного метода с базовой событийной UDP-синхронизацией без прикладной надежности.

	КЛЮЧЕВЫЕ СЛОВА

	согласованность состояний, многопользовательское приложение, распределенная асинхронная система, событийная синхронизация, конечный автомат, UDP, подтверждение доставки, повторная отправка, контроль порядка сообщений.
\end{essay}
